\documentclass[12pt]{article}
\usepackage{fullpage,enumerate,amsmath,bm,array,mathtools}
\usepackage[onehalfspacing]{setspace}

\DeclareMathOperator*{\argmax}{arg\,max}
\DeclareMathOperator*{\argmin}{arg\,min}
\DeclareMathOperator*{\tr}{tr}
\DeclareMathOperator*{\rk}{rk}
\DeclareMathOperator*{\spa}{span}
\DeclareMathOperator*{\ve}{vec}
\DeclareMathOperator*{\diag}{diag}
\DeclareMathOperator*{\Rea}{Re}
\DeclareMathOperator*{\var}{Var}
\DeclareMathOperator*{\cov}{Cov}

\begin{document}


\begin{center}
{\large ECMA 31150 -- Spring 2026} \\[0.5\baselineskip]
{\Large \bf Midterm Exam} \\[0.3\baselineskip]
April 21, 2026 \\[0.3\baselineskip]
Instructor: Mikkel Plagborg-Moller
\end{center}

\bigskip

\begin{itemize}
\item You have 75 minutes to complete the following 5 problems.
\item The test is closed-book. You may consult a calculator and one double-sided sheet of handwritten notes.
\item The total number of points is 100.
\item The problems have not been ordered in terms of difficulty. Many of the questions can be answered independently of each other.
\item \textbf{Show your work}. State several significant digits for every numerical answer.
\item \textbf{Partial answers will receive partial credit.} Do not leave questions empty. However, points may be subtracted for long, irrelevant remarks.
\item Good luck!
\end{itemize}

\clearpage

\subsection*{Problem 1 (10 points total)}
Let the time series $\lbrace y_t \rbrace$ be defined by
\[y_t \equiv (u_t + u_{t-1})u_{t-1} \quad \text{for } t=1,2,\dots,\quad \text{where} \quad u_t \overset{iid}{\sim} N(0,1) \quad \text{for } t=0,1,2,\dots\]
\begin{enumerate}[i)]
\item (10 points.) Show that $\lbrace y_t \rbrace$ is covariance stationary. Derive its mean, its ACF, and its long-run variance. [Hint: the 4th moment of the standard normal distribution equals 3.]
\end{enumerate}

\subsection*{Problem 2 (20 points total)}
Let $\lbrace y_t \rbrace$ be a covariance stationary process with autocovariance function $\gamma(\ell)\equiv \cov(y_t,y_{t-\ell})$. Let $\hat{\sigma}_y^2 \equiv \frac{1}{T-1}\sum_{t=1}^T (y_t-\bar{y})^2$ denote its sample variance, where $\bar{y} \equiv \frac{1}{T}\sum_{t=1}^T y_t$.

\begin{enumerate}[i)]
\item (12 points.) Prove that
\[E\left[\hat{\sigma}_y^2\right] = \gamma(0) - \frac{2}{T-1}\sum_{\ell=1}^{T-1} \left(1-\frac{\ell}{T}\right) \gamma(\ell).\]
[Hint: it is without loss of generality to assume $E[y_t]=0$.]
\item (8 points.) Give a concrete example of a covariance stationary model for $\lbrace y_t \rbrace$ such that $\hat{\sigma}_y^2$ is a strictly \emph{downward-biased} estimator of the population variance $\gamma(0)$.
\end{enumerate}

\subsection*{Problem 3 (20 points total)}
Consider the AR(2) process
\[y_t = -2 - \tfrac{1}{4}y_{t-1} + \tfrac{1}{8}y_{t-2} + u_t,\]
where $\lbrace u_t \rbrace$ is a white noise process with variance 2, and $E[u_t \mid \lbrace y_s \rbrace_{s<t}]=0$.

\begin{enumerate}[i)]
\item (8 points.) Verify that the process is covariance stationary, provided the initial values have appropriate moments. [Note: you do not need to derive the required moments of the initial values.]
\item (12 points.) Give an explicit formula for the optimal two-step-ahead forecast $y_{T+2|T} \equiv E[y_{T+2} \mid \lbrace y_s \rbrace_{s \leq T}]$? What is the associated RMSFE?
\end{enumerate}

\subsection*{Problem 4 (26 points total)}
We have $T=196$ observations of two time series $\lbrace y_t,x_t \rbrace$. We assume an ADL model for $y_t$ with two autoregressive lags and one distributed lag of $x_t$. A regression produces the following output:
\medskip

{\small \setstretch{0.9}
\begin{verbatim}
                            OLS Regression Results                            
==============================================================================
Dep. Variable:                      y   R-squared:                       0.469
Model:                            OLS   Adj. R-squared:                  0.460
Method:                 Least Squares   F-statistic:                     49.42
Date:                Tue, 21 Apr 2026   Prob (F-statistic):           1.18e-23
Time:                        11:00:00   Log-Likelihood:                -274.71
No. Observations:                 194   AIC:                             557.4
Df Residuals:                     190   BIC:                             570.5
Df Model:                           3                                         
Covariance Type:                  HC1                                         
==============================================================================
                 coef    std err          z      P>|z|      [0.025      0.975]
------------------------------------------------------------------------------
const         -0.0101      0.072     -0.140      0.889      -0.152       0.132
y_lag1         0.6422      0.070      9.172      0.000       0.505       0.779
y_lag2        -0.0719      0.062     -1.156      0.247      -0.194       0.050
x_lag1         0.1971      0.053      3.728      0.000       0.093       0.301
==============================================================================

Notes:
[1] Standard Errors are heteroscedasticity robust (HC1)
\end{verbatim}
}

\begin{enumerate}[i)]
\item (6 points.) Suppose the last observations in the data set are $y_T = 1$, $x_T = 0$, $y_{T-1}=-1$, $x_{T-1}=3$. What is the optimal forecast of $y_{T+1}$?
\item (6 points.) Propose a null hypothesis for the model coefficients under which $\lbrace y_t \rbrace$ is white noise. Then test this hypothesis, using a 5\% significance level. [Note: if you're unsure how to use the regression output to answer this question, you will receive full points if you provide a detailed general description of how to conduct the test.]
\item (6 points.) Test whether $\lbrace x_t \rbrace$ Granger-causes $\lbrace y_t \rbrace$, using a 10\% significance level.
\item (8 points.) Assume that $\lbrace y_t,x_t\rbrace$ are jointly covariance stationary. Treating the above coefficient estimates as exact rather than estimated, derive an expression for $E[y_t]$ as a function of $E[x_t]$.
\end{enumerate}

\subsection*{Problem 5 (24 points total)}
Indicate whether each of the following statements is \underline{true} or \underline{false}. Justify your answer with a very brief argument.

\begin{enumerate}[i)]
\item (8 points.) Suppose we observe a sample of size $T$ from the time series $\lbrace \tilde{y}_t \rbrace$ given by $\tilde{y}_t \equiv (1-\tfrac{1}{2}L)^{-1}y_t$, where $\lbrace y_t \rbrace$ follows a stationary AR(1) process with non-zero AR coefficient. Then the BIC procedure will select 2 lags with arbitrarily high probability as $T \to \infty$.

\item (8 points.) A bivariate VAR(1) process with coefficient matrix
\[\bm{A} = \begin{pmatrix}
0.8 & 0 \\
1.5 & -0.3
\end{pmatrix}\]
is covariance stationary (provided the initial value has appropriate moments).

\item (8 points.) If we run the recursive residual bootstrap for a VAR(1) model, using a very large number of iterations, then the sample average of the bootstrap draws $\hat{A}_{ij}^*$ will be very close to the real-data parameter estimate $\hat{A}_{ij}$. [Note: an informal argument suffices for this question.]
\end{enumerate}


\end{document}
 